\documentclass[12pt]{article}
\usepackage[utf8]{inputenc}
\usepackage{amsmath, amssymb, amsthm}
\usepackage{graphicx}
\usepackage{geometry}
\usepackage{fancyhdr}

\geometry{letterpaper, margin=1in}
\newcommand{\Mod}[1]{\ (\mathrm{mod}\ #1)}
\setlength{\headheight}{15pt}

% Header setup
\pagestyle{fancy}
\fancyhf{}
\lhead{Math 300 Spring 2026}
\rhead{Jeffrey Lu}
\cfoot{\thepage}
\parskip=1ex

\title{Exam 2 6\\ \large Math 300, Spring 2026 \\ Rutgers University}
\author{Jeffrey Lu}

\begin{document}

\maketitle

\noindent{}\textbf{Question 1:}

\begin{proof}
    By definition of even/odd, for some integers $a,b$, we have $x = 2a$ and $y = 2b+1$.
    Then,
    \begin{align*}
        z & = (2a)^2 + (2b+1)^2                       \\
          & = 4a^2 + 4b^2 + 4b + 1                    \\
          & = 4k +1, \text{ where }k = a^2 + b^2 + b.
    \end{align*}
    So $k$ is an integer. Then $z \in R$, and so $L \subseteq R$.

    Consider $k= -1$. Then, $4k+1 = -3 \in R$.
    But, $L$ is the sum of squares, so all elements in $L$ are positive.
    Therefore, $R \not\subseteq L$.
\end{proof}

\noindent{}\textbf{Question 2:}

(a)
\begin{proof}
    Let $S = (0,1)$.
    For sake of contradiction, assume that $S$ has a greatest element, say $x$.
    Since $x \in S$, $0 < x < 1$.

    Consider $y = \frac{x+1}{2} < \frac{1+1}{2} < 1$.
    Then, $y \in S$.
    But, $y > \frac{x+x}{2} = x$.
    This is a contradiction since $y > x$.
    Thus, $S$ does not have a largest element.
\end{proof}

(b)
\begin{proof}
    Assume that $A, B, C, D$ are sets satisfying $A \subseteq B$ and $C \subseteq D$.
    Suppose $(x,y) \in A \times C$.
    Then, $x \in A$ and $y \in C$.
\end{proof}

\noindent{}\textbf{Question 3:}
\begin{proof}
    Let $S = \{ 1, 3, 5, 7, 9, 11, 13, 15, 17, 19, 21, 23, 25, 27, 29 \}$.
    Then, we partition it into $P = \{
        \{3, 9, 27\},
        \{5, 15, 25 \},
        \{7, 21 \},
        \{1\},
        \{11 \},
        \{13\},
        \{17\},
        \{19\},
        \{23\},
        \{29\}
        \}$.
    Let these subsets of $S$ be the pigeonholes, and the 11 selected numbers be the pigeons.
    Then, by the pigeonhole principle, there 2 of the selected numbers must be in the same partition of $S$.
    Since all elements of the partitions of $S$ with more than 1 element are divisible by each other,
    we have found our pair.
\end{proof}



\noindent{}\textbf{Question 4:}
\noindent{}\textbf{Question 6:}



\end{document}

